\section{Обзор литературы и используемые модели}

\subsection{Машинно-обучаемые межатомные потенциалы}

Машинно-обучаемые межатомные потенциалы приближают потенциальную энергетическую поверхность атомной системы. Для конфигурации \(X\), состоящей из атомных типов, координат и, при необходимости, периодической ячейки, модель строит приближение энергии
\[
    \widehat{E}_{\theta}(X),
\]
а силы связываются с производными этой энергии по координатам атомов:
\[
    \widehat{\mathbf{F}}_{\theta,i}
    =
    -\nabla_{\mathbf{r}_i}\widehat{E}_{\theta}(X).
\]
Эта связь уже была использована в постановке задачи, но в обзоре моделей она задает главный критерий выбора архитектур: модель должна не только возвращать скалярную энергию, но и задавать достаточно гладкую и геометрически согласованную энергетическую поверхность, чтобы производные давали физически осмысленные силы.

Большинство современных машинно-обучаемых потенциалов используют локальное разложение энергии. В обобщенной форме его можно записать как
\[
    \widehat{E}_{\theta}(X)
    =
    \sum_{i=1}^{N} \widehat{E}_{\theta,i}(\mathcal{N}_i),
\]
где \(\mathcal{N}_i\) --- локальное окружение атома \(i\), определяемое соседями в радиусе отсечения. Это разложение не является специфичным для одной модели: оно задает общий язык для MTP, M3GNet, MACE и MatterSim. Различия между ними начинаются в том, как именно кодируется окружение \(\mathcal{N}_i\), как учитываются симметрии и используется ли предобучение.

Для данной работы важны две стратегии построения потенциала. Первая стратегия --- обучить специализированную модель только на целевых данных. Так используются MTP-потенциалы. Вторая стратегия --- взять универсальную предобученную модель и адаптировать ее к целевой системе. Так используются MatterSim и MACE. Поэтому обзор ниже не перечисляет все существующие MLIP, а разбирает модели, которые задают это сравнение.

\subsection{MTP как специализированный потенциал}

Moment Tensor Potentials были предложены как систематически улучшаемый класс межатомных потенциалов, в котором локальное окружение атома описывается через моментные тензоры \cite{shapeev2016mtp}. Для атома \(i\) с соседями \(j\) в радиусе отсечения вводятся относительные радиус-векторы
\[
    \mathbf{r}_{ij} = \mathbf{r}_j - \mathbf{r}_i.
\]
Моментный тензор можно записать в обобщенном виде как
\[
    M_{\mu,\nu}(\mathcal{N}_i)
    =
    \sum_{j \in \mathcal{N}_i}
    f_{\mu}\!\left(|\mathbf{r}_{ij}|, z_i, z_j\right)
    \mathbf{r}_{ij}^{\otimes \nu},
    \label{eq:mtp_moment_tensor}
\]
где \(z_i\) и \(z_j\) --- химические типы атомов, \(f_{\mu}\) --- радиальная функция, а \(\mathbf{r}_{ij}^{\otimes \nu}\) обозначает тензорное произведение радиус-вектора с самим собой \(\nu\) раз. Эта формула показывает, что MTP кодирует локальное окружение не только через расстояния, но и через направления на соседей. После свертки индексов из таких тензоров строятся инвариантные базисные функции \(B_{\alpha}(\mathcal{N}_i)\), а атомный вклад в энергию задается линейной комбинацией:
\[
    \widehat{E}_{\theta,i}(\mathcal{N}_i)
    =
    \sum_{\alpha}
    c_{\alpha} B_{\alpha}(\mathcal{N}_i).
    \label{eq:mtp_atomic_energy}
\]
Здесь обучаемыми параметрами являются коэффициенты \(c_{\alpha}\). Для этой работы это важно: MTP выступает как специализированный потенциал с явно контролируемым базисом, а не как большая предобученная нейросетевая модель.

Практическая работа с MTP в проекте выполнялась через MLIP-4 \cite{mlip4software}. Статью MLIP-3 можно использовать как контекст по современной линии MLIP-пакетов для построения MTP и активного обучения, но не как точное описание использованной версии MLIP-4 \cite{podryabinkin2023mlip3}. Поэтому в тексте разделяются три уровня: статья Shapeev описывает математическую основу MTP, MLIP-4 фиксирует использованный программный инструментарий, а MLIP-3 дает контекст развития пакетов для MTP.

Уровни MTP-16 и MTP-20 в работе задают разные наборы базисных функций. В использованных запусках MTP-16 соответствует 122 базисным функциям, MTP-20 --- 390. Число обучаемых параметров при этом не определяется только уровнем потенциала: оно также зависит от набора химических элементов и радиального базиса. В экспериментах использовался \texttt{RadialBasisCinf}\allowbreak\texttt{(size=8)}; при этих настройках число параметров составляло 608 и 932 для organic-датасета, 222 и 450 для H2O, 897 и 1293 для MoNbTaVW для уровней MTP-16 и MTP-20 соответственно. Эти числа нужны не для ранжирования моделей сами по себе, а для интерпретации MTP как специализированной базовой линии с контролируемой сложностью.

Роль MTP в сравнении не сводится к роли заведомо простой модели. Он обучается на тех же целевых данных, что и дообучаемые MatterSim и MACE, но не использует широкое предобучение. Поэтому сравнение с MTP показывает, насколько тонкая настройка универсального потенциала дает выигрыш относительно модели, которая строит приближение только из целевого набора данных.

\subsection{M3GNet как графовая модель атомистических систем}

M3GNet строит межатомный потенциал на графовом представлении структуры \cite{m3gnet2022}. Атомы задаются вершинами, связи между соседними атомами --- ребрами, координаты атомов и матрица ячейки входят в описание структуры. В статье это представлено как граф
\[
    G = (V, E, X, [M, u]),
\]
где \(V\) содержит атомные признаки, \(E\) --- признаки связей, \(X\) --- координаты атомов, \(M\) --- матрицу ячейки, а \(u\) --- дополнительные глобальные признаки, например температуру или давление, если они используются. Такая запись важна для нашей задачи: она явно отделяет химическую информацию от геометрической и позволяет затем получать силы через дифференцирование энергии по координатам.

Ключевой вклад M3GNet --- добавление трехчастичных взаимодействий в графовую модель. В обобщенном виде обновление ребра \(e_{ij}\) через \(n\)-частичное взаимодействие в статье записывается как
\[
    \widetilde{e}_{ij}
    =
    \sum_{\substack{
        k_1,k_2,\ldots,k_{n-2}\in \mathcal{N}_i / j \\
        k_1 \ne k_2 \ne \ldots \ne k_{n-2}
    }}
    \phi_n
    \left(
        e_{ij}, \mathbf{r}_{ij}, v_j,
        \mathbf{r}_{ik_1}, \ldots, \mathbf{r}_{ik_{n-2}},
        v_{k_1}, \ldots, v_{k_{n-2}}
    \right).
    \label{eq:m3gnet_n_body}
\]
Для M3GNet основное внимание уделяется случаю трехчастичных взаимодействий. Тогда информация о соседнем атоме \(k\) используется не только через расстояния, но и через угол \(\theta_{jik}\) между связями \(ij\) и \(ik\). В статье это реализуется через разложение углового взаимодействия по сферическим функциям и радиальному базису:
\[
    \widetilde{e}_{ij}
    =
    \sum_k
    j_l\!\left(z_{ln}\frac{r_{ik}}{r_c}\right)
    Y_l^0(\theta_{jik})
    \odot
    \sigma(W_v v_k + b_v)
    f_c(r_{ij})f_c(r_{ik}),
    \label{eq:m3gnet_three_body}
\]
где \(j_l\) --- сферическая функция Бесселя, \(Y_l^0\) --- сферическая гармоника, \(f_c\) --- функция отсечения, а \(\odot\) обозначает покомпонентное произведение. Для отчета не требуется подробно разбирать каждый множитель; важен смысл формулы: M3GNet добавляет в обновление связи информацию о тройках атомов и углах, то есть расширяет описание локальной геометрии за пределы попарных расстояний.

После учета многочастичных взаимодействий M3GNet выполняет последовательные обновления ребер, атомных признаков и, при наличии, глобального состояния:
\[
    e'_{ij}
    =
    e_{ij}
    +
    \phi_e(v_i \oplus v_j \oplus e_{ij} \oplus u)
    W_e^0 e_{ij}^0,
    \label{eq:m3gnet_edge_update}
\]
\[
    v'_i
    =
    v_i
    +
    \sum_j
    \phi'_e(v_i \oplus v_j \oplus e'_{ij} \oplus u)
    W_e^{0'} e_{ij}^0.
    \label{eq:m3gnet_atom_update}
\]
Здесь \(\oplus\) обозначает конкатенацию признаков. Эти формулы показывают, что информация накапливается через локальные связи, но атомный признак после нескольких обновлений уже содержит сведения о более широком окружении.

Для построения межатомного потенциала M3GNet предсказывает энергию, а силы и напряжения получает через автоматическое дифференцирование:
\[
    \mathbf{f}
    =
    -\frac{\partial E}{\partial \mathbf{x}},
    \qquad
    \sigma
    =
    V^{-1}\frac{\partial E}{\partial \varepsilon},
    \label{eq:m3gnet_forces_stress}
\]
где \(\mathbf{x}\) --- координаты атомов, \(V\) --- объем, \(\varepsilon\) --- деформация \cite{m3gnet2022}. Для нашего сравнения это принципиально: качество по силам проверяет не независимую векторную голову модели, а согласованность выученной энергетической поверхности с ее производными.

В статье M3GNet функция потерь включает энергию, силы и, для неорганических систем, напряжения:
\[
    \mathcal{L}
    =
    \ell(e, e_D)
    +
    w_f \ell(\mathbf{f}, \mathbf{f}_D)
    +
    w_{\sigma}\ell(\sigma,\sigma_D),
    \label{eq:m3gnet_loss}
\]
где \(\ell\) --- функция потерь Хьюбера, а индекс \(D\) обозначает данные из квантово-механического расчета \cite{m3gnet2022}. Эта формула важна как пример того, что обучение межатомного потенциала обычно балансирует несколько физических величин. В нашей работе итоговая оценка также разделяется на ошибки энергии и сил, но сравнение проводится уже через MAE и RMSE, введенные в разделе 2.

\subsection{MatterSim как универсальная предобученная модель}

MatterSim предложен как универсальная предобученная модель для атомистического моделирования в широком пространстве химических составов и термодинамических условий \cite{mattersim2024}. В статье модель позиционируется как потенциал, обученный на конфигурациях, покрывающих элементы периодической таблицы, температуры от \(0\) до \(5000\,\mathrm{K}\) и давления до \(1000\,\mathrm{GPa}\). Для этой работы важен не сам масштаб предобучения как рекламная характеристика, а его экспериментальная роль: MatterSim начинает адаптацию не со случайной инициализации, а с параметров, уже полученных на широком наборе атомистических данных.

Архитектурно MatterSim не нужно сводить к одной модели. В статье указано, что в MatterSim используются две графовые основы: M3GNet и Graphormer \cite{mattersim2024}. M3GNet описан выше как графовая модель с трехчастичными взаимодействиями. Graphormer относится к трансформерным графовым архитектурам; в MatterSim он используется для вариантов, где требуется более высокая емкость модели, но цена такой емкости --- более высокая стоимость вывода и требования к памяти. В этом отчете нельзя утверждать, что все использованные контрольные точки MatterSim имеют один и тот же backbone, если это не подтверждено логами конкретного запуска.

Обобщенно MatterSim можно описать как отображение
\[
    X
    \longmapsto
    \kappa_G
    \longmapsto
    \widehat{E}_{\theta}(X),
\]
где \(X\) --- атомная структура, а \(\kappa_G\) --- глобальное представление графа после графовых обновлений или структурного кодировщика. В дополнительных материалах MatterSim такое глобальное представление далее преобразуется через отдельный декодер. Для M3GNet-ветки это можно записать как
\[
    \kappa'_G
    =
    \mathrm{MLP}_1(\kappa_G)
    =
    \phi\!\left(
        W^2 \rho(W^1 \kappa_G + b_1) + b_2
    \right),
    \label{eq:mattersim_mlp_m3gnet}
\]
а для Graphormer-ветки как
\[
    \kappa'_G
    =
    \mathrm{MLP}_2(\kappa_G)
    =
    W^2 \lambda\!\left(
        \rho(W^1 \kappa_G + b_1)
    \right)
    + b_2.
    \label{eq:mattersim_mlp_graphormer}
\]
Эти формулы не нужны для ручного воспроизведения MatterSim, но они показывают, что модель разделяет две части: получение представления структуры и декодирование этого представления в целевую физическую величину.

Для тонкой настройки MatterSim используется обычная идея адаптации предобученной модели: параметры \(\theta_{\mathrm{pre}}\), полученные на широком наборе данных, становятся начальной точкой оптимизации на целевом наборе:
\[
    \theta_{\mathrm{ft}}
    =
    \arg\min_{\theta}
    \mathcal{L}_{D_{\mathrm{target}}}(\theta),
    \qquad
    \theta \text{ инициализируется из } \theta_{\mathrm{pre}}.
    \label{eq:mattersim_finetune}
\]
Документация MatterSim описывает такую настройку как штатный режим работы с моделью \cite{mattersimfinetuningdocs}. Для нашего сравнения это означает, что MatterSim представляет стратегию «широкое предобучение + адаптация к целевой системе». В разделе результатов эта стратегия сравнивается с MTP, который обучается на целевых данных без такого предобучения.

\subsection{MACE как эквивариантная модель для силовых полей}

MACE строится как графовая модель с эквивариантной передачей сообщений, то есть message passing \cite{mace2022}. В стандартной схеме графовой передачи сообщений состояние атома \(i\) на слое \(t\) можно записать как
\[
    \sigma_i^{(t)}
    =
    \left(\mathbf{r}_i, z_i, h_i^{(t)}\right),
    \label{eq:mace_node_state}
\]
где \(\mathbf{r}_i\) --- координаты, \(z_i\) --- химический элемент, \(h_i^{(t)}\) --- обучаемые признаки. Сообщение к атому строится через соседей:
\[
    m_i^{(t)}
    =
    \bigoplus_{j\in \mathcal{N}(i)}
    M_t\!\left(\sigma_i^{(t)},\sigma_j^{(t)}\right),
    \label{eq:mace_general_message}
\]
после чего признаки обновляются:
\[
    h_i^{(t+1)}
    =
    U_t\!\left(\sigma_i^{(t)},m_i^{(t)}\right).
    \label{eq:mace_general_update}
\]
В MACE, как и в других потенциалах с локальным разложением энергии, итоговая энергия собирается из атомных вкладов:
\[
    E_i
    =
    \sum_{t=1}^{T}
    R_t\!\left(\sigma_i^{(t)}\right).
    \label{eq:mace_readout_general}
\]

Эквивариантность означает, что внутренние признаки модели должны изменяться согласованно с поворотом или отражением атомной структуры. Для преобразования \(Q \in O(3)\) это условие записывается как
\[
    h_i^{(t)}\!\left(Q\cdot(\mathbf{r}_1,\ldots,\mathbf{r}_N)\right)
    =
    D(Q)\,
    h_i^{(t)}(\mathbf{r}_1,\ldots,\mathbf{r}_N),
    \label{eq:mace_equivariance}
\]
где \(D(Q)\) --- матрица представления преобразования \(Q\) \cite{mace2022}. Для скалярной энергии это особенно важно: при повороте всей структуры энергия не должна меняться, но векторные и тензорные внутренние признаки могут преобразовываться не тривиально. Такая конструкция дает модели возможность работать с геометрией, не нарушая физических симметрий.

Главная архитектурная идея MACE --- сообщения повышенного порядка. В статье сообщение записывается как иерархическое разложение по корреляционному порядку:
\[
    m_i^{(t)}
    =
    \sum_j u_1^{(t)}(\sigma_i^{(t)};\sigma_j^{(t)})
    +
    \sum_{j_1,j_2}
    u_2^{(t)}(\sigma_i^{(t)};\sigma_{j_1}^{(t)},\sigma_{j_2}^{(t)})
    + \cdots +
    \sum_{j_1,\ldots,j_{\nu}}
    u_{\nu}^{(t)}
    (\sigma_i^{(t)};\sigma_{j_1}^{(t)},\ldots,\sigma_{j_{\nu}}^{(t)}).
    \label{eq:mace_body_order_expansion}
\]
Эта формула показывает отличие MACE от моделей, где каждое сообщение фактически описывает только пару атомов. MACE пытается включить более сложные локальные взаимодействия непосредственно в сообщение, а не только накапливать их большим числом слоев.

Для построения таких сообщений MACE сначала формирует двухчастичные признаки \(A_i^{(t)}\), в которых используются радиальные функции, сферические гармоники и признаки соседнего атома. В упрощенной записи первого слоя это выглядит как
\[
    A_{i,klm}^{(1)}
    =
    \sum_{j\in \mathcal{N}(i)}
    R_{kl}^{(1)}(r_{ji})
    Y_{lm}(\widehat{\mathbf{r}}_{ji})
    W_{kz_j}^{(1)}.
    \label{eq:mace_a_features_simplified}
\]
Здесь \(R_{kl}^{(1)}\) кодирует расстояние, \(Y_{lm}\) --- направление на соседа, а \(W_{kz_j}^{(1)}\) зависит от химического типа атома \(j\). Затем признаки повышенного порядка строятся через тензорные произведения и симметризацию:
\[
    B_{i,\eta_{\nu} kLM}^{(t)}
    =
    \sum_{\mathbf{l}\mathbf{m}}
    C_{\eta_{\nu},\mathbf{l}\mathbf{m}}^{LM}
    \prod_{\xi=1}^{\nu}
    \sum_{\widetilde{k}}
    w_{k\widetilde{k}l_{\xi}}^{(t)}
    A_{i,\widetilde{k}l_{\xi}m_{\xi}}^{(t)}.
    \label{eq:mace_b_features}
\]
Коэффициенты \(C_{\eta_{\nu},\mathbf{l}\mathbf{m}}^{LM}\) обеспечивают нужные свойства преобразования признаков, а произведение по \(\xi\) соответствует корреляционному порядку. Для отчета важно не техническое устройство этих коэффициентов, а смысл: MACE строит признаки, которые кодируют многочастичную геометрию окружения и при этом сохраняют контролируемое поведение при поворотах и отражениях структуры.

Сообщение MACE после этого записывается как линейная комбинация таких признаков:
\[
    m_{i,kLM}^{(t)}
    =
    \sum_{\nu}
    \sum_{\eta_{\nu}}
    W_{z_i kL,\eta_{\nu}}^{(t)}
    B_{i,\eta_{\nu} kLM}^{(t)}.
    \label{eq:mace_message}
\]
Обновление признаков включает само сообщение и остаточную связь с предыдущим состоянием:
\[
    h_{i,kLM}^{(t+1)}
    =
    \sum_{\widetilde{k}}
    W_{kL,\widetilde{k}}^{(t)}
    m_{i,\widetilde{k}LM}^{(t)}
    +
    \sum_{\widetilde{k}}
    W_{z_i kL,\widetilde{k}}^{(t)}
    h_{i,\widetilde{k}LM}^{(t)}.
    \label{eq:mace_update}
\]
Энергия считывается только из инвариантных признаков \(h_{i,k00}^{(t)}\), чтобы атомные вклады в энергию оставались скалярными:
\[
    E_i
    =
    E_i^{(0)} + E_i^{(1)} + \cdots + E_i^{(T)},
    \qquad
    E_i^{(t)}
    =
    R_t\!\left(\{h_{i,k00}^{(t)}\}_k\right).
    \label{eq:mace_energy_readout}
\]
Эта часть архитектуры напрямую связана с задачей регрессии силовых полей: модель использует геометрически богатые внутренние признаки, но итоговая энергия сохраняет правильную инвариантность.

В данной работе MACE представляет вторую универсальную стратегию наряду с MatterSim: предобученная модель дообучается на целевой системе и затем сравнивается с MTP. Документация MACE описывает тонкую настройку как отдельный рабочий режим \cite{macefinetuningdocs}. При этом MACE default и MACE-EW50 должны рассматриваться как разные протоколы настройки, а не как одна и та же модельная строка. Это различение будет использовано в разделе результатов для MoNbTaVW.

\subsection{Что именно сравнивается в этой работе}

Рассмотренные модели различаются не только реализациями, но и источником информации, доступной до обучения на целевой системе. MTP строит приближение с нуля по данным конкретной системы. MatterSim и MACE начинают с предобученных параметров и затем адаптируются на целевом наборе. Поэтому основное сравнение в работе можно записать как сравнение двух отображений:
\[
    D_{\mathrm{target}}
    \longmapsto
    \widehat{E}_{\mathrm{MTP}},
    \qquad
    (\theta_{\mathrm{pre}}, D_{\mathrm{target}})
    \longmapsto
    \widehat{E}_{\mathrm{universal,ft}},
    \label{eq:comparison_strategies}
\]
где \(D_{\mathrm{target}}\) --- целевой набор данных, \(\theta_{\mathrm{pre}}\) --- параметры предобученной универсальной модели, а \(\widehat{E}_{\mathrm{universal,ft}}\) --- энергия после тонкой настройки MatterSim или MACE.

Эта запись подчеркивает, что сравниваются не только архитектуры, но и режимы использования данных. MTP проверяет, чего можно добиться, обучаясь на целевой системе без внешнего предобучения. MatterSim и MACE проверяют, насколько помогает перенос информации из широкого предобучения при последующей адаптации. Такой дизайн особенно полезен для трех разных систем проекта: organic-датасета, H2O и MoNbTaVW, поскольку они различаются химическим составом и характером разбиений.

Из обзора также следует ограничение на интерпретацию результатов. Ни одна из архитектурных идей сама по себе не гарантирует лучшую ошибку по всем метрикам. MTP имеет контролируемый специализированный базис, M3GNet и MatterSim используют графовое представление атомистических данных, MACE дополнительно встраивает эквивариантность и сообщения повышенного порядка. Поэтому дальнейшие выводы должны формулироваться не как общее ранжирование моделей, а как утверждения вида: система, разбиение, модель, метрика.

В разделе результатов это означает, что отдельно анализируются Energy MAE, Energy RMSE, Forces MAE и Forces RMSE; отдельно различаются organic-датасет, H2O и MoNbTaVW; отдельно фиксируются MACE default и MACE-EW50. Такая дисциплина нужна, чтобы сравнение оставалось проверяемым: универсальная модель может быть сильнее в одном режиме и не быть лучшей в другом, а изменение протокола тонкой настройки может менять баланс ошибок энергии и сил.
